% Chapter 12 converted from Markdown
% Auto-generated - do not edit manually




\section*{Section 12.1: Before PLIx: Evidence Chains}

Before PLIx, SEG (Shared Evidence Graph) stores evidence chains—linking claims to evidence (code, docs, tests, decisions) but lacking intent awareness.

\textbf{SEG's Original Purpose}

SEG was designed to:

\begin{itemize}
\item \textbf{Store Evidence Chains:} Link claims to evidence through graph edges
\item \textbf{Track Entities:} Store entities (claims, sources, derivations, agents)
\item \textbf{Track Relations:} Store relations (SUPPORTS, CONTRADICTS, REFERENCES)
\item \textbf{Enable Reasoning:} Enable queries like "what evidence supports this claim?"
\end{itemize}

SEG's strength lies in its graph-based structure, enabling complex evidence reasoning through entity-relation graphs.

\textbf{Evidence Chain Example}

Before PLIx, SEG stores evidence chains:

\begin{lstlisting}[language=Python, caption=Python Example]
# Create claim entity
claim = Entity(
    type="claim",
    name="Room booking system works correctly",
    attributes={"description": "System can book rooms"}
)

claim_entity = seg.add_entity(claim)

# Create evidence entity
evidence = Entity(
    type="evidence",
    name="Test results",
    attributes={"test_file": "test_booking.py", "pass_rate": 0.95}
)

evidence_entity = seg.add_entity(evidence)

# Create evidence relation
relation = Relation(
    source_id=evidence_entity.id,
    target_id=claim_entity.id,
    relation_type=RelationType.SUPPORTS,
    confidence=0.95
)

seg.add_relation(relation)

# Query: What evidence supports this claim?
supporting_evidence = seg.query_relations(
    target_id=claim_entity.id,
    relation_type=RelationType.SUPPORTS
)
\end{lstlisting}

SEG stores evidence chains (what supports what) but doesn't track intent lineage (what intents led to outcomes). This limits SEG's ability to reason about purpose and verify intent-outcome relationships.

\textbf{Limitations of Evidence Chains}

Evidence chains have limitations:

\begin{itemize}
\item \textbf{No Intent Awareness:} SEG doesn't know what was intended, only what evidence exists
\item \textbf{No Intent Lineage:} Can't trace outcomes back to intents
\item \textbf{No Intent Evolution:} Can't track how intent evolved over time
\item \textbf{No Intent-Outcome Mapping:} Can't map outcomes to intents
\end{itemize}

These limitations prevent SEG from supporting intent-driven reasoning, verification, and learning.

\textbf{Entity-Relation Structure}

SEG uses entity-relation structure:

\begin{lstlisting}[language=Python, caption=Python Example]
# Entities represent claims, sources, derivations, agents
claim_entity = Entity(type="claim", name="Room booking works")
source_entity = Entity(type="source", name="Test results")
agent_entity = Entity(type="agent", name="Test runner")

# Relations link entities
support_relation = Relation(
    source_id=source_entity.id,
    target_id=claim_entity.id,
    relation_type=RelationType.SUPPORTS
)
\end{lstlisting}

Entity-relation structure enables complex reasoning, but without intent awareness, entities don't represent intents.

\textbf{Before PLIx Summary}

Before PLIx, SEG is execution-focused:
\begin{itemize}
\item Stores evidence chains (what supports what)
\item Tracks entities and relations
\item Enables evidence reasoning
\item Lacks intent awareness (no intent lineage or evolution tracking)
\end{itemize}

This execution focus limits SEG's ability to support intent-driven reasoning and verification.

\section*{Section 12.2: After PLIx: Intent Lineage}

After PLIx, SEG stores intent lineage—tracing outcomes back to intents, tracking intent evolution, and mapping intent-outcome relationships.

\textbf{Intent-Aware Entities}

With PLIx, SEG stores intent entities:

\begin{lstlisting}[language=Python, caption=Python Example]
# Create intent entity
intent_entity = Entity(
    type="intent",
    name="Book a meeting room",
    attributes={
        "contract": contract.to_dict(),
        "intent_type": "booking",
        "domain": "meeting_rooms"
    }
)

intent_entity_id = seg.add_entity(intent_entity)

# Create outcome entity
outcome_entity = Entity(
    type="outcome",
    name="Room reserved",
    attributes={
        "room_reserved": True,
        "reservation_id": "res-123"
    }
)

outcome_entity_id = seg.add_entity(outcome_entity)

# Create intent-outcome relation
intent_outcome_relation = Relation(
    source_id=intent_entity_id,
    target_id=outcome_entity_id,
    relation_type=RelationType.ACHIEVES,  # Intent \textrightarrow{} Outcome
    confidence=0.90,
    attributes={
        "postconditions_satisfied": True,
        "verification_timestamp": datetime.now()
    }
)

seg.add_relation(intent_outcome_relation)
\end{lstlisting}

SEG now stores intent entities and intent-outcome relations, enabling intent lineage tracking.

\textbf{Intent Lineage Tracking}

With PLIx, SEG tracks intent lineage:

\begin{lstlisting}[language=Python, caption=Python Example]
# Track intent lineage: NL \textrightarrow{} Contract \textrightarrow{} Plan \textrightarrow{} Execution \textrightarrow{} Outcome
nl_intent_entity = Entity(
    type="nl_intent",
    name="Book a meeting room",
    attributes={"original_text": "Book a meeting room"}
)

contract_entity = Entity(
    type="plix_contract",
    name="PLIx Contract",
    attributes={"contract": contract.to_dict()}
)

plan_entity = Entity(
    type="execution_plan",
    name="APOE Execution Plan",
    attributes={"plan": plan.to_dict()}
)

execution_entity = Entity(
    type="execution",
    name="Execution Result",
    attributes={"result": result.to_dict()}
)

outcome_entity = Entity(
    type="outcome",
    name="Room Reserved",
    attributes={"room_reserved": True}
)

# Create lineage chain
seg.add_relation(Relation(
    source_id=nl_intent_entity.id,
    target_id=contract_entity.id,
    relation_type=RelationType.COMPILES_TO
))

seg.add_relation(Relation(
    source_id=contract_entity.id,
    target_id=plan_entity.id,
    relation_type=RelationType.COMPILES_TO
))

seg.add_relation(Relation(
    source_id=plan_entity.id,
    target_id=execution_entity.id,
    relation_type=RelationType.EXECUTES_TO
))

seg.add_relation(Relation(
    source_id=execution_entity.id,
    target_id=outcome_entity.id,
    relation_type=RelationType.PRODUCES
))

# Query lineage: Trace outcome back to NL intent
lineage = seg.query_lineage(outcome_entity.id, direction="backward")
# Returns: Outcome \textrightarrow{} Execution \textrightarrow{} Plan \textrightarrow{} Contract \textrightarrow{} NL Intent
\end{lstlisting}

Intent lineage tracking enables SEG to trace outcomes back to intents, enabling queries like "what intent led to this outcome?"

\textbf{Intent Evolution Tracking}

With PLIx, SEG tracks intent evolution:

\begin{lstlisting}[language=Python, caption=Python Example]
# Store intent version 1
intent_v1 = Entity(
    type="intent",
    name="Book a meeting room",
    attributes={"version": 1, "contract": contract_v1.to_dict()}
)

intent_v1_id = seg.add_entity(intent_v1)

# Store intent version 2 (evolved)
intent_v2 = Entity(
    type="intent",
    name="Book a meeting room with catering",
    attributes={"version": 2, "contract": contract_v2.to_dict()}
)

intent_v2_id = seg.add_entity(intent_v2)

# Create evolution relation
evolution_relation = Relation(
    source_id=intent_v1_id,
    target_id=intent_v2_id,
    relation_type=RelationType.EVOLVES_TO,
    attributes={
        "evolution_type": "refinement",
        "changes": ["added_catering_requirement"]
    }
)

seg.add_relation(evolution_relation)

# Query evolution: How did intent evolve?
evolution_chain = seg.query_lineage(intent_v2_id, relation_type=RelationType.EVOLVES_TO)
# Returns: Intent v1 \textrightarrow{} Intent v2 (evolution chain)
\end{lstlisting}

Intent evolution tracking enables SEG to track how intent evolved over time, enabling queries like "how did this intent evolve?"

\textbf{Intent-Outcome Mapping}

With PLIx, SEG maps outcomes to intents:

\begin{lstlisting}[language=Python, caption=Python Example]
# Map outcome to intent
def map_outcome_to_intent(outcome: dict, seg: SEGraph) -> List[Entity]:
    """Map outcome to intents that achieved it"""
    
    # Find outcomes matching this outcome
    outcome_entities = seg.query_entities(
        type="outcome",
        attributes_filter={"room_reserved": True}
    )
    
    # Find intents that achieved these outcomes
    intent_entities = []
    for outcome_entity in outcome_entities:
        relations = seg.query_relations(
            target_id=outcome_entity.id,
            relation_type=RelationType.ACHIEVES
        )
        for relation in relations:
            intent_entity = seg.get_entity(relation.source_id)
            intent_entities.append(intent_entity)
    
    return intent_entities

# Query: What intents achieved this outcome?
intents = map_outcome_to_intent({"room_reserved": True}, seg)
\end{lstlisting}

Intent-outcome mapping enables SEG to query which intents achieved which outcomes, enabling intent-driven learning.

\textbf{After PLIx Summary}

After PLIx, SEG is intent-aware:
\begin{itemize}
\item Stores intent lineage (traces outcomes back to intents)
\item Tracks intent evolution (how intent evolved over time)
\item Maps intent-outcome relationships (which intents achieved which outcomes)
\item Enables intent-driven reasoning (queries about intent and outcomes)
\end{itemize}

This intent awareness transforms SEG from execution-focused evidence to intent-aware evidence, enabling intent-driven reasoning and learning.

\section*{Section 12.3: Transformation Details}

The transformation from evidence chains to intent lineage involves storing PLIx contracts as SEG entities, creating intent relations, collecting intent evidence, and enabling intent lineage queries.

\textbf{PLIx Contracts \textrightarrow{} SEG Entities}

PLIx contracts store as SEG entities:

\begin{lstlisting}[language=Python, caption=Python Example]
def store_plix_contract_as_entity(contract: PLIxContract, seg: SEGraph) -> str:
    """Store PLIx contract as SEG entity"""
    
    entity = Entity(
        type="plix_contract",
        name=contract.intent,
        attributes={
            "intent": contract.intent,
            "contract": contract.to_dict(),
            "tasks": [task.to_dict() for task in contract.tasks],
            "constraints": contract.constraints,
            "evidence": contract.evidence
        }
    )
    
    entity_id = seg.add_entity(entity)
    return entity_id
\end{lstlisting}

This transformation preserves contract semantics while enabling SEG's graph-based reasoning.

\textbf{Intent Relations}

Intent relations link intents to outcomes:

\begin{lstlisting}[language=Python, caption=Python Example]
def create_intent_outcome_relation(
    intent_entity_id: str,
    outcome_entity_id: str,
    verification_result: bool,
    confidence: float,
    seg: SEGraph
) -> str:
    """Create intent-outcome relation"""
    
    relation = Relation(
        source_id=intent_entity_id,
        target_id=outcome_entity_id,
        relation_type=RelationType.ACHIEVES,
        confidence=confidence,
        attributes={
            "postconditions_satisfied": verification_result,
            "verification_timestamp": datetime.now()
        }
    )
    
    relation_id = seg.add_relation(relation)
    return relation_id
\end{lstlisting}

Intent relations enable SEG to track which intents achieved which outcomes, enabling intent-outcome reasoning.

\textbf{Intent Evidence Collection}

Intent evidence collection stores evidence in SEG:

\begin{lstlisting}[language=Python, caption=Python Example]
def collect_intent_evidence(
    contract: PLIxContract,
    outcome: dict,
    execution_provenance: dict,
    seg: SEGraph
) -> str:
    """Collect intent evidence and store in SEG"""
    
    # Create evidence entity
    evidence_entity = Entity(
        type="intent_evidence",
        name=f"Evidence for {contract.intent}",
        attributes={
            "contract": contract.to_dict(),
            "outcome": outcome,
            "execution_provenance": execution_provenance,
            "postconditions_satisfied": verify_contract(contract, outcome)
        }
    )
    
    evidence_id = seg.add_entity(evidence_entity)
    
    # Link to intent
    seg.add_relation(Relation(
        source_id=evidence_id,
        target_id=get_intent_entity_id(contract, seg),
        relation_type=RelationType.PROVIDES_EVIDENCE_FOR
    ))
    
    return evidence_id
\end{lstlisting}

Intent evidence collection enables SEG to store proof of intent achievement, supporting verification and learning.

\textbf{Intent Lineage Queries}

Intent lineage queries enable intent-driven reasoning:

\begin{lstlisting}[language=Python, caption=Python Example]
def query_intent_lineage(outcome_entity_id: str, seg: SEGraph) -> List[Entity]:
    """Query intent lineage: Trace outcome back to intent"""
    
    # Find relations where outcome is target
    relations = seg.query_relations(
        target_id=outcome_entity_id,
        relation_type=RelationType.ACHIEVES
    )
    
    # Get intent entities
    intent_entities = []
    for relation in relations:
        intent_entity = seg.get_entity(relation.source_id)
        if intent_entity.type == "plix_contract":
            intent_entities.append(intent_entity)
    
    return intent_entities

def query_outcome_lineage(intent_entity_id: str, seg: SEGraph) -> List[Entity]:
    """Query outcome lineage: Trace intent to outcomes"""
    
    # Find relations where intent is source
    relations = seg.query_relations(
        source_id=intent_entity_id,
        relation_type=RelationType.ACHIEVES
    )
    
    # Get outcome entities
    outcome_entities = []
    for relation in relations:
        outcome_entity = seg.get_entity(relation.target_id)
        outcome_entities.append(outcome_entity)
    
    return outcome_entities
\end{lstlisting}

Intent lineage queries enable SEG to trace outcomes to intents and intents to outcomes, enabling intent-driven reasoning.

\textbf{Transformation Benefits}

The transformation provides:

\begin{itemize}
\item \textbf{Intent Lineage:} SEG tracks intent lineage, enabling outcome-to-intent tracing
\item \textbf{Intent Evolution:} SEG tracks intent evolution, enabling temporal reasoning
\item \textbf{Intent-Outcome Mapping:} SEG maps outcomes to intents, enabling learning
\item \textbf{Intent Evidence:} SEG stores intent evidence, enabling verification
\end{itemize}

These benefits transform SEG from execution-focused evidence to intent-aware evidence, enabling intent-driven reasoning and learning.

\section*{Section 12.4: Implementation Examples}

Implementation examples demonstrate PLIx \textrightarrow{} SEG entity creation, intent relation creation, intent evidence collection, and intent lineage queries.

\textbf{Example 1: PLIx \textrightarrow{} SEG Entity Creation}

\begin{lstlisting}[language=Python, caption=Python Example]
# PLIx contract
contract = PLIxContract(
    intent="Book a meeting room",
    contract={"post": ["room_reserved == true"]}
)

# Store as SEG entity
intent_entity_id = store_plix_contract_as_entity(contract, seg)

# Query entity
intent_entity = seg.get_entity(intent_entity_id)
print(f"Intent: {intent_entity.attributes['intent']}")
print(f"Contract: {intent_entity.attributes['contract']}")
\end{lstlisting}

This example demonstrates storing PLIx contracts as SEG entities, enabling graph-based reasoning.

\textbf{Example 2: Intent Relation Creation}

\begin{lstlisting}[language=Python, caption=Python Example]
# Create intent-outcome relation
intent_entity_id = store_plix_contract_as_entity(contract, seg)
outcome_entity_id = seg.add_entity(Entity(
    type="outcome",
    name="Room Reserved",
    attributes={"room_reserved": True}
))

relation_id = create_intent_outcome_relation(
    intent_entity_id,
    outcome_entity_id,
    verification_result=True,
    confidence=0.90,
    seg
)

print(f"Intent-outcome relation created: {relation_id}")
\end{lstlisting}

This example demonstrates creating intent-outcome relations, linking intents to outcomes.

\textbf{Example 3: Intent Evidence Collection}

\begin{lstlisting}[language=Python, caption=Python Example]
# Collect intent evidence
evidence_id = collect_intent_evidence(
    contract=contract,
    outcome={"room_reserved": True},
    execution_provenance={"execution_id": "exec-123"},
    seg=seg
)

print(f"Evidence collected: {evidence_id}")

# Query evidence
evidence_entity = seg.get_entity(evidence_id)
print(f"Postconditions satisfied: {evidence_entity.attributes['postconditions_satisfied']}")
\end{lstlisting}

This example demonstrates collecting intent evidence and storing it in SEG.

\textbf{Example 4: Intent Lineage Queries}

\begin{lstlisting}[language=Python, caption=Python Example]
# Query: What intents led to this outcome?
outcome_entity_id = seg.add_entity(Entity(
    type="outcome",
    name="Room Reserved",
    attributes={"room_reserved": True}
))

intents = query_intent_lineage(outcome_entity_id, seg)
print(f"Intents that achieved this outcome: {len(intents)}")
for intent in intents:
    print(f"  - {intent.attributes['intent']}")

# Query: What outcomes did this intent achieve?
intent_entity_id = store_plix_contract_as_entity(contract, seg)
outcomes = query_outcome_lineage(intent_entity_id, seg)
print(f"Outcomes achieved by this intent: {len(outcomes)}")
\end{lstlisting}

This example demonstrates intent lineage queries, tracing outcomes to intents and intents to outcomes.

\textbf{Implementation Benefits}

Implementation examples demonstrate:

\begin{itemize}
\item \textbf{Entity Creation:} Storing PLIx contracts as SEG entities
\item \textbf{Relation Creation:} Creating intent-outcome relations
\item \textbf{Evidence Collection:} Collecting and storing intent evidence
\item \textbf{Lineage Queries:} Querying intent lineage for reasoning
\end{itemize}

These examples show how SEG transforms from execution-focused evidence to intent-aware evidence, enabling intent-driven reasoning and learning.

\section*{Chapter 12 Summary}

SEG transforms from evidence chains to intent lineage through PLIx integration. Before PLIx, SEG stores evidence chains but lacks intent awareness. After PLIx, SEG stores intent lineage, tracks intent evolution, maps intent-outcome relationships, and enables intent-driven reasoning. This transformation enables intent-driven verification, learning, and reasoning, making SEG a foundation for intent-aware systems.

\textbf{Next:} Part III Integration complete. Part IV explores implementation—CNL compiler, runtime, adapters, and testing.

\textbf{Word Count:} \textasciitilde{}2,300 words  

